% Einheit 1 – Monotonie und erste Ableitung
\setcounter{aufgabe}{8}
\einheitenkopf[e1][1 Monotonie]{Einheit 1 von 5 · Monotonie und erste Ableitung}
\zweigzeile{Hier lernst du, am Vorzeichen der Ableitung abzulesen und am Term nachzuweisen, wo ein Graph steigt oder fällt · kennst du seit Q1 · Abitur GK · baut auf: Gleichungen lösen, Ableitungen bilden (Nr. 2–5)}

\verfahren{Monotonie aus der Ableitung bestimmen}
\begin{aufgabe}{Ich kann am Vorzeichen der Ableitung ablesen, ob der Graph steigt oder fällt.}
Gib an, ob $f'(x)$ im Intervall positiv oder negativ ist und ob $f$ dort steigt oder fällt.
\begin{teile}
\teil $f'(x) = x - 1$ auf $[2;\,5]$: \leerfeld
\teil $f'(x) = x - 1$ auf $[-3;\,0]$: \leerfeld
\teil $f'(x) = -2x + 8$ auf $[5;\,9]$: \leerfeld
\teil $f'(x) = (x - 1)(x - 5)$ auf $[2;\,4]$: \leerfeld
\end{teile}
\end{aufgabe}

\begin{aufgabe}{Ich kann die Monotonieintervalle einer Funktion bestimmen.}
Bestimme, wo $f$ steigt und wo $f$ fällt. Gib die Intervalle an.
\begin{teile}
\teil $f(x) = x^3 - 12x + 1$
\teil $f(x) = -x^3 + 3x^2 + 9x$
\anweisung{Gib ohne Rechnung an, wo $f$ steigt und wo $f$ fällt.}
\teil $f$ ist ganzrational dritten Grades und hat den Hochpunkt $H(-1\,|\,5)$ und den Tiefpunkt $T(3\,|\,{-}1)$.
\anweisung{Rechne mit der Ableitung.}
\teil Gegeben sind $f(x) = x^3 - 5x^2 + 9x + 1$ und $g(x) = x^2 + 1$. Die Funktion $d$ mit $d(x) = f(x) - g(x)$ gibt für $0 \le x \le 5$ den senkrechten Abstand der beiden Graphen an. Bestimme, wo $d$ auf $[0;\,5]$ steigt und wo $d$ fällt, und gib den Wertebereich von $d$ auf $[0;\,5]$ an. (Abitur 2021 GK)
\end{teile}
\end{aufgabe}

\verfahren{Monotonie am Term nachweisen}
\begin{aufgabe}{Ich kann am Term der Ableitung nachweisen, dass ein Graph überall steigt.}
Zeige, dass $f$ auf ganz $\mathbb{R}$ streng monoton steigt.
\begin{teilezwei}
\tz $f(x) = 2x^3 + x$ & \tz $f(x) = 4 + e^{2x}$ \\
\tz $f(x) = x + e^{x}$ & \tz $f(x) = 5 - e^{-x}$ \\
\tz $f(x) = x^3 + x^2 + x$ & \\
\anweisung{Begründe, dass der Graph von $f$ keinen Hochpunkt und keinen Tiefpunkt hat.}
\tz $f(x) = x^5 + 2x$ & \\
\anweisung{Weise die Monotonie am Term nach.}
\tz $f(x) = (x^2 + 3)\cdot e^{x}$ \quad (Abitur 2024 GK) & \\
\end{teilezwei}
\end{aufgabe}

\begin{aufgabe}{Ich finde den Fehler bei Monotonieintervallen.}
Finde den Fehler, benenne ihn und korrigiere die Rechnung.
\begin{teile}
\teil Tim untersucht $f(x) = x^3 - 6x^2 + 5$:
\rechnung{f''(x) &= 6x - 12 \\ f''(x) &> 0 \quad \text{für } x > 2}
Tim schreibt: „$f$ fällt für $x \le 2$ und steigt für $x \ge 2$.“
\teil Mia untersucht $f(x) = x^2 - 6x + 1$:
\rechnung{f'(x) &= 2x - 6}
Mia schreibt: „$f$ fällt für $x \le 0$ und steigt für $x \ge 0$.“
\end{teile}
\end{aufgabe}

\begin{aufgabe}{Ich kann begründen, ob ein Graph überall steigt.}
Entscheide ohne Rechnung und begründe.
\begin{teile}
\teil Warum wird $4x^2 + 1$ für kein $x$ null oder negativ?
\teil Für $f(x) = x^3 + 2$ gilt $f'(0) = 0$. Jonas sagt: „Dann steigt $f$ nicht auf ganz $\mathbb{R}$ streng monoton.“ Hat er recht?
\end{teile}
\end{aufgabe}

\begin{aufgabe}{Ich kann die Monotoniebereiche zweier Modelle vergleichen.}
Die Temperatur in einem Gewächshaus zwischen 6 und 16 Uhr wird durch zwei Modelle beschrieben ($t$ in Stunden nach 6 Uhr, $0 \le t \le 10$, Temperatur in $^\circ$C):
\[ A(t) = -0{,}1t^3 + 1{,}2t^2 + 12 \qquad\text{und}\qquad B(t) = 12 + 6t\cdot e^{-0{,}2t}. \]
\begin{teile}
\teil Bestimme für beide Modelle, in welchem Zeitraum die Temperatur steigt.
\teil Beurteile die Aussage: „Nach Modell A steigt die Temperatur drei Stunden länger als nach Modell B.“ (IQB 2023 grundlegend)
\end{teile}
\end{aufgabe}
